\documentclass[10pt,a4paper]{amsart}
\usepackage[margin=19mm]{geometry}
\usepackage{amsmath,amssymb,mathtools}
\usepackage[scr=rsfs,cal=euler]{mathalfa}
\usepackage{xfrac}
\usepackage[unicode,hidelinks,bookmarks=false]{hyperref}
\newcommand{\ie}{i.e.\ }
\newcommand{\cf}{cf.\ }
\newcommand{\resp}{respectively\ }
\newcommand{\Subsection}{\subsection}
\providecommand{\nobreakdash}{\nobreak\hbox{-}}
\setlength{\emergencystretch}{2em}
\multlinegap0pt
\usepackage{amssymb}
\newtheorem{theorem}{Theorem}
\newtheorem{lemma}{Lemma}
\def\T{ \mathbb{T} }
\def\b{ \mathbf{B} }
\def\div{ \hbox{\rm div}\,  }
\newcommand{\norm}[2]{\left\lVert #1 \right\rVert_{#2}}
\def\u{ \mathbf{u} }
\title{Stabilization by a background magnetic field: global well-posedness of the compressible isentropic ideal MHD equations with velocity damping}
\author{Liening Qiao, Jiahong Wu, Fuyi Xu$^{\dag}$, Xiaoping Zhai}
\date{Source: arXiv:2605.04462v3 (CC BY 4.0)}
\begin{document}
\maketitle

\noindent\emph{Source and licence.} Stabilization by a background magnetic field: global well-posedness of the compressible isentropic ideal MHD equations with velocity damping. Version arXiv:2605.04462v3 (CC BY 4.0), distributed by arXiv under the Creative Commons Attribution 4.0 International licence, http://creativecommons.org/licenses/by/4.0/, as recorded in the arXiv licence metadata for this exact version and shown on the abstract page https://arxiv.org/abs/2605.04462v3. Full licence notice: \texttt{licenses/arxiv-2605.04462-CC-BY-4.0.txt}.\par\smallskip

\noindent\textit{Editorial scope.} Counted unit: the article's theorem dingli (source0.tex lines 516--539). The article's complete original proof of it is delivered with it (source0.tex lines 1185--1349, one \texttt{proof} environment of 165 lines, which the article places away from the statement). No mathematics is written, repaired or re-derived: the statement and the proof are the article's own lines, in the article's own notation, and no retained line is substituted or re-flowed.  Retained uncounted as the source-local context the proof needs: lines 246--263; lines 433--437; lines 727--736; lines 750--760; lines 785--787; lines 828--840; lines 961--965; lines 978--993; lines 1010--1026; lines 1143--1153.  Nothing was omitted from any retained span, so the package carries no omission record.  The retained lines cite 3 works, whose entries are transcribed verbatim from the article's own bibliography source in the e-print and delivered with short editorial labels recorded in the item's reference mapping.  No retained line contains a figure environment, an image-inclusion command or drawing code, so no image is delivered and the package declares no asset dependency.  Editorial formatting only, in addition: an AMS \texttt{amsart} preamble that restates the article's own declarations and macros used by the retained lines, namely the environments \texttt{theorem}, \texttt{lemma} on separate counters, together with the standard packages those lines need.  The retained environments print in this document as Theorem 1, Lemma 1 in order of appearance, whereas the article prints the same items with the numbering of its own full document.  The complete native source of the article as distributed in the arXiv e-print for this exact version is bundled unchanged as \texttt{sources/199862/source0.tex}.
\begin{equation}\label{MHD1}
	\left\{
	\begin{aligned}
		&\partial_t a + \mathrm{div}(\rho\mathbf{u}) = 0, \\[4pt]
		&\rho(\partial_t\mathbf{u} + \mathbf{u}\cdot\nabla\mathbf{u}) + \rho\mathbf{u}
		+ \nabla P
		= \underbrace{\omega\cdot\nabla\b
			- \nabla(\omega\cdot\b)}_{\text{linear in }\b}
		+ \underbrace{\b\cdot\nabla\b
			- \tfrac{1}{2}\nabla|\b|^2}_{\text{nonlinear}}, \\[4pt]
		&\partial_t\b + \mathbf{u}\cdot\nabla\b
		= \underbrace{\omega\cdot\nabla\mathbf{u} - \omega\,\mathrm{div}\,\mathbf{u}}_{\text{linear in }\mathbf{u}}
		+ \underbrace{\b\cdot\nabla\mathbf{u} - \b\,\mathrm{div}\,\mathbf{u}}_{\text{nonlinear}}, \\[4pt]
		&\mathrm{div}\,\b = 0, \\[4pt]
		&(\rho,\mathbf{u},\b)|_{t=0} = (\rho_0,\mathbf{u}_0,\b_0).
	\end{aligned}
	\right.
\end{equation}

\begin{equation}\label{Diophantine0}
	\exists\; c>0 \;\text{ and }\; r>2 \;\text{ such that }\;
	|\omega\cdot\mathbf{k}| \geq \frac{c}{|\mathbf{k}|^r}, \qquad
	\forall\;\mathbf{k}\in\mathbb{Z}^3\setminus\{\mathbf{0}\}.
\end{equation}

\begin{lemma}[Kato--Ponce product estimate, {\cite{kato}}]
	\label{Hk-estimate}
	Let $s\geq 0$.  For any
	$f,g\in H^s(\mathbb{T}^3)\cap L^\infty(\mathbb{T}^3)$,
	\begin{equation}\label{KP-product}
		\|fg\|_{H^s}
		\leq C\bigl(\|f\|_{L^\infty}\|g\|_{H^s}
		+ \|g\|_{L^\infty}\|f\|_{H^s}\bigr).
	\end{equation}
\end{lemma}

\begin{lemma}[Sobolev estimate for compositions, {\cite{Triebel}}]
	\label{Hk-estimate-1}
	Let $s>0$, $f\in H^s(\mathbb{T}^3)\cap L^\infty(\mathbb{T}^3)$, and
	let $F\in C^\infty(\mathbb{R})$ with $F(0)=0$.  Then
	\begin{equation}\label{composition}
		\|F(f)\|_{H^s}
		\leq C\bigl(1+\|f\|_{L^\infty}\bigr)^{[s]+1}\|f\|_{H^s},
	\end{equation}
	where the constant $C$ depends on
	$\sup_{k\leq[s]+2,\,|t|\leq\|f\|_{L^\infty}}|F^{(k)}(t)|$.
\end{lemma}

\begin{equation}\begin{split}\label{pri-2}
		\frac{1}{2}\le \rho(t,x) \le \frac32, \quad \text{for all } (x,t)\in\mathbb T^3\times[0,T].
\end{split}\end{equation}

\begin{lemma}\label{lemma-pri1}
Let $(\rho,\u,\b) \in C([0, T];H^N)$ be a solution to the  system \eqref{MHD1}. Denote
$$E_s(t):=\int\Big(\rho|\u|^2+|\b|^2+e(\rho)\Big)\,dx+\Big\|\big(\frac{p'}{\rho}\Lambda^sa,\sqrt{\rho}\Lambda^s\u,\Lambda^s\b\big)\Big\|^2_0.$$ Then, for any $3\leq s\leq N$, it holds that
\begin{align}\label{pri-2-2}
\frac{d}{dt}E_s+
\|(u,\Lambda^s\u)\|_0^2\leq C(1+\|(a,\u,\b)\|_3)\|(a,\u,\b)\|_3\|(a,\u,\b)\|^2_s+
\|a\|_3\|\b\|^{2}_3\|a\|_s\|\u\|_s
\end{align}
with
\begin{equation}
\begin{split}\label{bu-2026-6-27-1}& E_s\geq\min\{\inf_{\rho\in(\frac12,\frac32)}\frac{p'(\rho)}{\rho},\frac12\}\big(\|(a,\u,\b)\|^2_0+\|\Lambda^s(a,\u,\b)\|^2_0\big),\\&E_s\leq\max\{\sup_{\rho\in(\frac12,\frac32)}\frac{p'(\rho)}{\rho},\frac32\}\big(\|(a,\u,\b)\|^2_0+\|\Lambda^s(a,\u,\b)\|^2_0\big),\end{split}\end{equation}
where the positive constant $C$ depends only on $|\omega|,\beta,$ and $Q$.
\end{lemma}

\begin{equation}\begin{split}\label{pri-5}
\|\partial_t\u\|_{L^\infty}&=\Big\|-(\u\cdot\nabla)\u-\u-\frac{p'}{\rho}\nabla a+\rho^{-1}\big((\omega\cdot \nabla) \b+(\b\cdot \nabla) \b-\nabla(\omega\cdot \b)-\frac12\nabla(|\b|^2)\big)\Big\|_{L^\infty}\\
&\lesssim\|\u\|_{L^\infty}\|\nabla\u\|_{L^\infty}+\|\u\|_{L^\infty}+\|\nabla a\|_{L^\infty}+(1+\|\b\|_{L^\infty})\|\nabla\b\|_{L^\infty}\\
&\lesssim\big(1+\|(\u,\b)\|_{2}\big)\|(a,\u,\b)\|_3.
\end{split}\end{equation}

\begin{equation}\begin{split}\label{pri-6}
\|\partial_t\u\|_{s-1}&=\left\|-\u\cdot\nabla \u-\u-\frac{p'}{\rho}\nabla a+\rho^{-1}\big((\omega\cdot \nabla) \b+(\b\cdot \nabla) \b-\nabla(\omega\cdot \b)-\frac12\nabla(|\b|^2)\big)\right\|_{s-1}\\
&\lesssim\|\u\|_{L^\infty}\|\nabla\u\|_{s-1}+\|\u\|_{s-1}\|\nabla\u\|_{L^\infty}+\|\u\|_{s-1}+\Big\|\frac{p'}{\rho}\Big\|_{s-1}\|\nabla a\|_{L^\infty}\\
&\quad+\Big\|\frac{p'}{\rho}\Big\|_{L^\infty}\|\nabla a\|_{s-1}+\Big\|\frac1\rho\Big\|_{L^\infty}
\|(\omega\cdot \nabla) \b+(\b\cdot \nabla) \b-\nabla(\omega\cdot \b)-\frac12\nabla(|\b|^2)\|_{s-1}\\
&\quad+\Big\|\frac1\rho\Big\|_{s-1}
\|(\omega\cdot \nabla) \b+(\b\cdot \nabla) \b-\nabla(\omega\cdot \b)-\frac12\nabla(|\b|^2)\|_{L^\infty}\\
&\lesssim\|\u\|_3\|\u\|_s+(1+\|a\|_s)\|a\|_3+\|(a,\u)\|_s+(1+\|\b\|_{L^\infty})\|\b\|_{s}\\
&\quad+(1+\|a\|_{s-1})\|\nabla \b\|_{L^\infty}
+(1+\|a\|_s)\|\nabla \b\|_{L^\infty}\|\b\|_{L^\infty}\\
&\lesssim\|\u\|_3\|\u\|_s+(1+\|a\|_s)\|a\|_3+\|(a,\u)\|_s+(1+\|\b\|_2)\|\b\|_{s}\\
&\quad+(1+\|a\|_{s})\|\b\|_3
+(1+\|a\|_s)\|\b\|_3^{2}
\\ &\lesssim\|\u\|_3\|\u\|_s+\|a\|_3\|a\|_s+\|(a,\u,\b)\|_s+\|\b\|_2\|\b\|_{s}\\
&\quad+\|\b\|^{2}_3+\|\b\|_3\|a\|_s+\|a\|_s\|\b\|^{2}_3.
\end{split}\end{equation}

\begin{equation}\label{MHD2}
\left\{
\begin{array}{rl}
&\hspace{-0,5cm}\partial_ta+\div \u=f_1,
\smallskip\\
&\hspace{-0,5cm}\partial_t \u
+\beta\nabla a +\u=(\omega\cdot \nabla) \b-\nabla(\omega\cdot  \b)+f_2,
\medskip\\
&\hspace{-0,5cm}\partial_t  \b
=(\omega\cdot \nabla)\u-\omega\div \u+f_3,
\medskip\\
&\hspace{-0,5cm}{\mathop{\rm div}}\, \b=0,
\medskip\\
&\hspace{-0,5cm}(\rho,\u, \b)_{|t=0}=(\rho_0,\u_0, \b_0),
\end{array}
\right.
\end{equation}

\begin{lemma}\label{lemma-pri2}Suppose $\omega$ satisfies the Diophantine condition \eqref{Diophantine0}
with constants $c$ and $r$. Let $(a,\u,\b) \in C([0, T];H^N)$ be a solution to the  system \eqref{MHD2}.
Then, it holds for any $t\in[0,T]$, $s\geq r$, that
\begin{equation}
\begin{split}\label{pri-199-9LLLL}
&c_1\|(a,\b)\|_{s-1}^2+\frac{d}{dt}\int \Big(\omega\cdot\nabla\Lambda^{s}a\cdot\Lambda^{s}(\u\cdot\omega)-\Lambda^{s}\u\cdot(\omega\cdot\nabla)\Lambda^{s}\b\Big) \,dx\\
&\quad\lesssim \|\u\|_{s+1}^{2}+\|(f_1,f_2,f_3)\|_{s}^{2}
\end{split}
\end{equation}
with a positive constant $c_1$ depending only on $\beta,c,r,|\omega|$ and $|\mathbb T^3|.$
\end{lemma}

\begin{theorem}\label{dingli}
	Let $N>3r+3$ with $r>2$, and assume $\omega\in\mathbb{R}^3$ satisfy
	the Diophantine condition \eqref{Diophantine0} for some constant $c>0$.
	Suppose the initial data $(\rho_0-1, \mathbf{u}_0, \b_0)\in H^N(\mathbb{T}^3)$
	satisfies the mean conditions
	\begin{equation}\label{junzhi0}
		\int_{\mathbb{T}^3}\rho_0\,dx = 1, \qquad
		\int_{\mathbb{T}^3}\b_0\,dx = 0,
	\end{equation}
	and $\frac12\leq\rho_0\leq \frac32.$  Then there exists $\varepsilon>0$ such that if
	\begin{equation}\label{2026-6-27-3-1}
	\|\rho_0-1\|_{H^N} + \|\mathbf{u}_0\|_{H^N}
	+ \|\b_0\|_{H^N} \leq \varepsilon,
	\end{equation}
	the system \eqref{MHD1} admits a unique global solution
	$(\rho-1, \mathbf{u}, \b)\in C([0,\infty); H^N)$.
	Moreover, for any $t\geq 0$ and $r+1\leq\gamma< N$,
	\begin{equation}\label{decay}
		\|\rho(t)-1\|_{H^{\gamma}}
		+ \|\mathbf{u}(t)\|_{H^\gamma}
		+ \|\b(t)\|_{H^\gamma}
		\leq C(1+t)^{-\frac{N-\gamma}{2r+2}}.
	\end{equation}
\end{theorem}

\begin{proof}[Proof of Theorem \ref{dingli}]  First of all, the MHD system in \eqref{MHD1} with any
initial data $(a_0, \u_0,  \b_0)$ in $ H^N(\T^3)$ has a unique local solution.
This follows from a standard contraction mapping argument (see, e.g.,
\cite{MM}). The
bootstrapping argument is employed to prove the global existence and stability.
It starts with the ansatz that the solution $(a(t), \u(t),
\b(t))$ to \eqref{MHD1} satisfies
\begin{align}\label{ping23}
\sup_{t\in[0,T]}(\|a\|_{{N}}+\|\u\|_{{N}}+\|\b\|_{{N}})\le \delta,
\end{align}
for some $0<\delta<1$ to be determined later. We then show that
\begin{align*}%\label{ping231}
\sup_{t\in[0,T]}(\|a\|_{{N}}+\|\u\|_{{N}}+\|\b\|_{{N}})\le \frac{\delta}{2}.
\end{align*}
The a priori assumption \eqref{pri-2} is guaranteed by the embedding relation and the smallness of the norm stated in \eqref{ping23}.
On the one hand, taking $s=r+1$ in \eqref{pri-2-2}, yields that
\begin{equation}\label{pri-2-2-22}
\begin{split}
\frac{d}{dt}E_{r+1}+\|(\u,\Lambda^{r+1}\u)\|_{0}^2
%&\quad\lesssim\|a\|_3\|\u\|_{4r+2}\big(\|\u\|_3\|\u\|_{4r+2}+\|a\|_3\|a\|_{4r+2}+\|(a,\b)\|_3+\|(a,\b)\|_{4r+2}\\
%&\qquad+\|\b\|_2\|\b\|_{4r+2}+\|\b\|^{2}_3+\|\b\|_3\|a\|_{4r+2}+\|a\|_{4r+2}\|\b\|^{2}_3\big)\\
%&\qquad+(1+\|(a,\u,\b)\|_3)\|(a,\u,\b)\|_3\|(a,\u,\b)\|^2_{4r+2}\\
\lesssim(1+\|(a,\u,\b)\|^2_3)\|(a,\u,\b)\|_3\|(a,\u,\b)\|^2_{r+1}.
\end{split}\end{equation}
On the other hand,  choosing  $s=r$ in \eqref{pri-199-9LLLL} implies that
\begin{equation}
\begin{split}\label{pri-199-9LLLL-L}
&c_1\|(a,\b)\|_{r-1}^2+\frac{d}{dt}\int \Big(\omega\cdot\nabla\Lambda^{r}a\cdot\Lambda^{r}(\u\cdot\omega)-\Lambda^{r}\u\cdot(\omega\cdot\nabla)\Lambda^{r}\b\Big) \,dx\\&\quad\leq C\|\u\|_{r+1}^{2}+C\|(f_1,f_2,f_3)\|_{r}^{2}.
\end{split}
\end{equation}
Hence, let
 $A$ be a positive  constant determined later. We infer from \eqref{pri-199-9LLLL-L} and \eqref{pri-2-2-22} that
\begin{equation}\label{pri-2-2-22-2A}
\begin{split}
&\frac{d}{dt}\left\{AE_{r+1}(t)+\int \Big(\omega\cdot\nabla\Lambda^{r}a\cdot\Lambda^{r}(\u\cdot\omega)-\Lambda^{r}\u\cdot(\omega\cdot\nabla)\Lambda^{r}\b\Big) \,dx\right\}\\
&\quad\quad+A\|(\u,\Lambda^{r+1}\u)\|_{0}^2-C\|\u\|_{r+1}^2+c_1\|(a,\b)\|_{r-1}^2\\
&\quad\lesssim (1+\|(a,\u,\b)\|^2_3)\|(a,\u,\b)\|_3\|(a,\u,\b)\|^2_{r+1}+\|(f_1,f_2,f_3)\|_{r}^{2}.
\end{split}\end{equation}
Define
\begin{align*}%\label{ping31}
{\mathcal{E}(t)}=&AE_{r+1}(t)+\int \Big(\omega\cdot\nabla\Lambda^{r}a\cdot\Lambda^{r}(\u\cdot\omega)-\Lambda^{r}\u\cdot(\omega\cdot\nabla)\Lambda^{r}\b\Big) \,dx.
\end{align*}
Then, from \eqref{bu-2026-6-27-1} and the
Cauchy--Schwarz inequality,  we may choose $A$ sufficiently large, depending only on $r$ and $\omega$ such that
\begin{equation}\label{bu2026-6-28-1}c_1\|\u\|_{r+1}^2\leq A\|(\u,\Lambda^{r+1}\u)\|_{0}^2-C\|\u\|^2_{r+1}
\end{equation} and
\begin{equation}\label{bu2026-6-28-2}
\|(a,\u,\b)\|^2_{r+1}\leq {\mathcal{E}(t)}\leq 2A\|(a,\u,\b)\|^2_{r+1}.
\end{equation}
We infer from \eqref{pri-2-2-22-2A} and \eqref{bu2026-6-28-1} that
\begin{equation}\label{bu-pri-2-2-22-2A}
\begin{split}
&\frac{d}{dt}{\mathcal{E}(t)}+c_1\|(a,\u,\b)\|_{0}^2\lesssim (1+\|(a,\u,\b)\|^2_3)\|(a,\u,\b)\|_3\|(a,\u,\b)\|^2_{r+1}+\|(f_1,f_2,f_3)\|_{r}^{2}.
\end{split}\end{equation}
It follows from Lemma \ref{Hk-estimate} and $\|\cdot\|_{L^\infty}\lesssim\|\cdot\|_2$ that
\begin{equation}\begin{split}\label{4.6AAAA}
\|f_1\|_{r}\lesssim\|a\u\|_{r+1}\lesssim\|(a,\u)\|_{L^\infty}\|(a,\u)\|_{r+1}\lesssim\|(a,\u)\|_{2}\|(a,\u)\|_{r+1},
\end{split}\end{equation}
and
\begin{equation}\begin{split}\label{bu4.7AAAA}
\|f_3\|_{r}\lesssim&\|(\u\cdot\nabla) \b\|_{r}+\|( \b\cdot\nabla)\u\|_{r}+\| \b\div \u\|_{r}\\
\lesssim&\|(\u,\nabla \b)\|_{L^\infty}\|(\u,\nabla  \b)\|_{r}+\| (\b,\nabla \u)\|_{L^\infty}\| (\b,\nabla \u)\|_r\\
\lesssim&\|(\u, \b)\|_3\|(\u, \b)\|_{r+1}.
\end{split}\end{equation}
It remains to bound the term $f_{2}.$  We  deduces by \eqref{pri-5}, \eqref{pri-6}, Lemma \ref{Hk-estimate},
and the embedding relation, that
\begin{equation*}
\begin{split}
\|a\partial_t\u\|_{r}&\lesssim\|a\|_{L^\infty}\|\partial_t\u\|_{r}+\|a\|_{r}\|\partial_t\u\|_{L^\infty}\\
&\lesssim \|a\|_3\big(\|\u\|_3\|\u\|_{r+1}+\|a\|_3\|a\|_{r+1}+\|(a,\u,\b)\|_{r+1}+\|\b\|_2\|\b\|_{r+1}\\
&\quad+\|\b\|^{2}_3+\|\b\|_3\|a\|_{r+1}+\|a\|_{r+1}\|\b\|^{2}_3\big)+\|a\|_{r}\big(1+\|\u,\b\|_{2}\big)\|(a,\u,\b)\|_3\\
&\lesssim (1+\|(a,\u,\b)\|_3)\|(a,\u,\b)\|_3\|(a,\u,\b)\|_{r+1}.
\end{split}
\end{equation*}
It follows from \eqref{pri-2}, Lemma \ref{Hk-estimate-1}, and the embedding inequality that
\begin{equation*}\begin{split}
\|(p'(a+1)-p'(1))\nabla a\|_{r}
&\lesssim\|p'(a+1)-p'(1)\|_{L^\infty}\|a\|_{r+1}+\|p'(a+1)-p'(1)\|_{r}\|\nabla a\|_{L^\infty}\\
&\lesssim\|a\|_{L^\infty}\|a\|_{r+1}+\|a\|_{r}\|\nabla a\|_{L^\infty}
\\&\lesssim\|a\|_3\|a\|_{r+1}.
\end{split}\end{equation*}
Employing  Lemma \ref{Hk-estimate} and \eqref{pri-2} yields that
\begin{equation*}\begin{split}
&\left\|-\rho(\u\cdot\nabla)\u-a\u+(\b\cdot\nabla)\b-\frac{1}{2}\nabla|\b|^2\right\|_{r}\\
&\quad\lesssim \|\rho \u,\nabla\u\|_{L^\infty}\|\rho \u,\nabla\u\|_{r}+\|(a,\u)\|_{L^\infty}\|(a,\u)\|_{r}+\|(\b,\nabla\b)\|_{L^\infty}\|(\b,\nabla\b)\|_{r}\\
&\quad\lesssim(\|\rho\|_{L^\infty}\|\u\|_{L^\infty}+\|\nabla\u\|_{L^\infty})(\|\rho\|_{L^\infty}\|\u\|_{r}+\|\rho\|_{r}\|\u\|_{L^\infty}
+\|\u\|_{r+1})
+\|(a,\u,\b)\|_3\|(a,\u,\b)\|_{r+1}\\
&\quad\lesssim\|\u\|_3(\|\u\|_{r+1}+\|\rho\|_{r}\|\u\|_{3})+\|(a,\u,\b)\|_3\|(a,\u,\b)\|_{r+1}\\
&\quad\lesssim\|(a,\u,\b)\|_3\|(a,\u,\b)\|_{r+1}+\|\u\|_3^2(1+\|a\|_{r})
\\&\quad\lesssim (1+\|(a,\u,\b)\|_3)\|(a,\u,\b)\|_3\|(a,\u,\b)\|_{r+1}.%\lesssim&(1+\|(a,\u),h\|_3)\|(a,\u),h\|_3\|(a,\u),h\|_{m+1}.
\end{split}\end{equation*}
Combining  with the above three inequalities,  we have
\begin{equation}\begin{split}\label{4.8AAAA}
\|f_2\|_{r}\lesssim (1+\|(a,\u,\b)\|_3)\|(a,\u,\b)\|_3\|(a,\u,\b)\|_{r+1}.
\end{split}\end{equation}
Then inserting \eqref{4.6AAAA}-\eqref{4.8AAAA} into \eqref{bu-pri-2-2-22-2A} gives rise to
\begin{equation}\label{pri-2-2-22-2}
\begin{split}
\frac{d}{dt}{\mathcal{E}(t)}+c_1\|(a,\u,\b)\|_{0}^2\lesssim \big(1+\|(a,\u,\b)\|_3\big)^3\|(a,\u,\b)\|_3\|(a,\u,\b)\|^2_{r+1}.
\end{split}\end{equation}
For any ${N}>3r+3$, by the interpolation inequality and the embedding relation, it holds that
\begin{equation*}\begin{split}
\|(a,\u,\b)\|^3_{r+1}&\lesssim \|(a,\u,\b)\|_{0}^{\frac{3(N-r-1)}{N}}\|(a,\u,\b)\|_{N}^{\frac{3(r+1)}{N}}
\\&\lesssim\|(a,\u,\b)\|_{0}^{2}\|(a,\u,\b)\|_{0}^{\frac{3(N-r-1)}{N}-2}\|(a,\u,\b)\|_{N}^{\frac{3(r+1)}{N}}
\\&\lesssim\|(a,\u,\b)\|_{0}^{2}\|(a,\u,\b)\|_{N}^{\frac{3(N-r-1)}{N}-2}\|(a,\u,\b)\|_{N}^{\frac{3(r+1)}{N}}
\\&\lesssim\|(a,\u,\b)\|_{0}^{2}\|(a,\u,\b)\|_{N},
\end{split}\end{equation*}
which together with \eqref{ping23} and the embedding relation yields that
\begin{equation*}\label{pri-2-2-22-4}
\begin{split}\big(1+\|(a,\u,\b)\|_3\big)^3\|(a,\u,\b)\|_3\|(a,\u,\b)\|^2_{r+1}
&\lesssim \big(1+\|(a,\u,\b)\|_3\big)^3\|(a,\u,\b)\|^3_{r+1}
\\&\lesssim\big(1+\|(a,\u,\b)\|_3\big)^3\|(a,\u,\b)\|_{N}\|(a,\u,\b)\|_{0}^{2}
\\&\lesssim\delta\|(a,\u,\b)\|_{0}^{2}.
\end{split}\end{equation*}
Plugging the above estimate  into \eqref{pri-2-2-22-2}, and then taking $\delta$ small enough such that $C\delta\leq\frac{c_1}{2}$,  we  deduce that
\begin{equation}\label{pri-2-2-22-6}
\frac{d}{dt}{\mathcal{E}(t)}+\frac{c_1}{2}\|(a,\u,\b)\|_{0}^2\leq0.
\end{equation}
For any ${N}> 3r+3$, by \eqref{bu2026-6-28-2}, the interpolation inequality and the  assumption (\ref{ping23}),  we have
\begin{equation*}%\label{pri-2-2-22-7}
 \big(\mathcal{E}(t)\big)^{\frac{N}{N-r-1}}\leq \big(2A\|(a,\u,\b)\|^2_{r+1}\big)^{\frac{N}{N-r-1}}\lesssim\|(a,\u,\b)\|_{0}^{2}\|(a,\u,\b)\|_{N}^{\frac{2r+2}{N-r-1}}
\leq C\delta^{\frac{2r+2}{N-r-1}}\|(a,\u,\b)\|_{0}^{2},
\end{equation*}
which along with \eqref{pri-2-2-22-6} yields the following Lyapunov-type inequality
\begin{equation*}%\label{pri-2-2-22-8-1}
\frac{d}{dt}{\mathcal{E}(t)}+c\mathcal{E}(t)^{\frac{N}{N-r-1}}\leq0.
\end{equation*}
It then follows easily that
\begin{align}\label{pri-2-2-22-8}
{\mathcal{E}(t)}\le C(1+t)^{-\frac{N-r-1}{r+1}}.
\end{align}
Taking $s={N}$ in \eqref{pri-2-2} and using $\eqref{bu-2026-6-27-1}_1$ and the fact, $\|\cdot\|^2_{N}\approx\|\cdot\|^2_{0}+\|\Lambda^N(\cdot)\|^2_{0},$ give rise to
\begin{equation}\label{pri-2-2-22-9}
\begin{split}
\frac{d}{dt}E_{N}+c_2\|\u\|_{N}^2&\leq C(1+\|(a,\u,\b)\|^2_3)\|(a,\u,\b)\|_3\|(a,\u,\b)\|^2_{N}\\
&\leq C(1+\|(a,\u,\b)\|^2_3)\|(a,\u,\b)\|_3E_{N}.
\end{split}\end{equation}
From \eqref{pri-2-2-22-8}, \eqref{bu2026-6-28-2}, and $N>3r+3$, we have
\begin{align*}
\int_0^T(1+\|(a,\u,\b)\|^2_3)\|(a,\u,\b)\|_3\,d\tau\lesssim\int_0^T(1+t)^{-\frac{N-r-1}{2r+2}}\,dt\le C.
\end{align*}
Then, applying Gronwall's inequality to \eqref{pri-2-2-22-9} and using \eqref{bu-2026-6-27-1} and \eqref{2026-6-27-3-1} yield
\begin{equation}\begin{split}\label{bu-pri-2-2-22-10}
\big\|(a,\u,\b)\big\|^2_{N}+\int_0^t\|\u(\tau)\|_{N}^2\,d\tau&\lesssim E_{N}(t)+c_2\int_0^T\|\u(\tau)\|_{N}^2d\tau\\
&\lesssim E_{N}(0)\lesssim\big\|(a(0),\u(0),\b(0))\big\|^2_{N}\le C\varepsilon^2.\end{split}\end{equation}
Furthermore, taking $\varepsilon$ small enough in \eqref{2026-6-27-3-1} so that $C\varepsilon^2\le \frac{\delta^2}{4}$, we deduce
\begin{align}\label{ping23-3}
\sup_{t\in[0,T]}(\|a\|_{{N}}+\|\u\|_{{N}}+\|\b\|_{{N}})\le \frac{\delta}{2},
\end{align}
which together with  the continuity argument  implies that the local solution
can be extended as a global one in time.
%Moreover, from \eqref{pri-2-2-22-8}, we also have the following decay rate
%\begin{align*}%\label{ping40}\norm{a(t)}{{4r+2}}+\norm{\u(t)}{{4r+2}}+\norm{\b(t)}{{4r+2}}\le  C(1+t)^{-\frac{N-4r-2}{2(2r+2)}}.\end{align*}

Moreover, for any $r+1\leq\gamma< N $,  thanks to the following interpolation inequality
\begin{align*}%\label{ming22}
	\norm{\cdot}{{\gamma}}\le\norm{\cdot}{{r+1}}^{\frac{{N}-\gamma}{{N}-r-1}}
	\norm{\cdot}{{N}}^{\frac{\gamma-r-1}{{N}-r-1}},
\end{align*}
we can conclude by \eqref{bu2026-6-28-2}, \eqref{pri-2-2-22-8} and \eqref{bu-pri-2-2-22-10} that  the   decay rate for the higher order energy
\begin{align*}%\label{ping44}
\norm{a(t)}{{\gamma}}+\norm{\u(t)}{{\gamma}}+\norm{\b(t)}{{\gamma}}\le C(1+t)^{-\frac{N-\gamma}{2r+2}}.
\end{align*}
This completes the proof of Theorem \ref{dingli}.\end{proof}

\begin{thebibliography}{99}
\bibitem[MM]{MM} A. Majda and A. Bertozzi, {\it Vorticity and Incompressible Flow}, Cambridge University Press, Cambridge, UK, 2002.
\bibitem[Triebe]{Triebel} H. Triebel, Theory of Function Spaces, Monogr. Math., Birkh$\ddot{\mathrm{a}}$user Verlag, Basel, Boston, 1983.
\bibitem[kato]{kato} T. Kato, Liapunov Functions and Monotonicity in the Euler and Navier-Stokes Equations, Lecture Notes in Mathematics, vol. 1450. Springer, Berlin (1990).
\end{thebibliography}
\end{document}
